\chapter{Fixed-Point Exclusion and Exceptional-Orbit Divisibility}
\label{chap:g017-fixed-point-exclusion}

\status{SOH-G017 proved: neither fixed point of the quotient negative inversion is a root of $F$; exceptional paired sets, if nonempty, have forced orbit divisibility.}

\section{The remaining fixed-point question}

Chapter~\ref{chap:quotient-negative-inversion-zero-set-no-go} introduced the quotient involution
\[
 J(w)=\frac{1}{16w}
\]
for the even entire function
\[
 \xi\!\left(\frac12+z\right)=F(z^2).
\]
Its two fixed values are $w=\pm1/4$.  The positive one is excluded by the previously proved positivity
\[
 F(x)>0\qquad(x\ge0).
\]
The only unresolved fixed candidate was therefore $w=-1/4$.

Because $z^2=-1/4$ gives $z=\pm i/2$,
\[
 F\!\left(-\frac14\right)
 =\xi\!\left(\frac12+\frac{i}{2}\right)
 =\Xi\!\left(\frac12\right),
\]
where $\Xi(t):=\xi(1/2+it)$.

\section{Positive block from the exact Riemann kernel}

From Chapter~\ref{chap:positive-riemann-kernel},
\[
 \Xi(t)=\int_0^\infty \Phi(y)\cos(ty)\,dy,
\]
with
\[
 \Phi(y)=\sum_{n\ge1}\phi_n(y),
\]
\[
 \phi_n(y)=4\pi n^2e^{5y/2}
 \bigl(2\pi n^2e^{2y}-3\bigr)e^{-\pi n^2e^{2y}}>0
 \qquad(y\ge0).
\]
At $t=1/2$,
\[
 F\!\left(-\frac14\right)
 =\int_0^\infty\Phi(y)\cos(y/2)\,dy.
\]

On $0\le y\le1/2$, retain only the $n=1$ channel.  Then
\[
 e^{5y/2}\ge1,
 \qquad
 2\pi e^{2y}-3\ge2\pi-3,
 \qquad
 e^{-\pi e^{2y}}\ge e^{-\pi e},
\]
and
\[
 \cos(y/2)\ge\cos(1/4).
\]
Hence
\[
 \int_0^{1/2}\Phi(y)\cos(y/2)\,dy\ge A,
\]
where
\[
 A:=2\pi(2\pi-3)e^{-\pi e}\cos(1/4).
\]
Using only $3<\pi<4$, $2<e<3$, and
\[
 \cos x\ge1-\frac{x^2}{2},
\]
we obtain
\[
 A>18\cdot\frac{31}{32}\cdot3^{-12}.
\]
The exact integer inequality
\[
 18\cdot31\cdot2^{15}>32\cdot3^{12}
\]
therefore gives
\begin{equation}
 \boxed{A>2^{-15}}.
 \label{eq:g017-positive-block}
\end{equation}

\section{Super-exponential oscillatory tail}

For $y\ge\pi$, set $r_n=\pi n^2e^{2y}$.  Since $2r_n-3<2r_n$,
\[
 \phi_n(y)
 <8\pi^2n^4e^{9y/2}e^{-\pi n^2e^{2y}}.
\]
Put $x=e^{2y}$ and $X=e^{2\pi}$.  Then
\[
 \int_\pi^\infty\phi_n(y)\,dy
 <4\pi^2n^4\int_X^\infty x^{5/4}e^{-\pi n^2x}\,dx.
\]
For $x\ge X$,
\[
 \frac{d}{dx}\log\!\left(x^{5/4}e^{-\pi n^2x}\right)
 =\frac{5}{4x}-\pi n^2
 \le-\frac{\pi n^2}{2},
\]
so
\[
 \int_X^\infty x^{5/4}e^{-\pi n^2x}\,dx
 \le\frac{2}{\pi n^2}X^{5/4}e^{-\pi n^2X}.
\]
Consequently
\[
 \int_\pi^\infty\Phi(y)\,dy\le T,
\]
with
\[
 T:=8\pi X^{5/4}\sum_{n\ge1}n^2e^{-\pi Xn^2}.
\]
If $q=e^{-\pi X}$, then $e^{-\pi Xn^2}\le q^n$ and
\[
 \sum_{n\ge1}n^2q^n=\frac{q(1+q)}{(1-q)^3}.
\]
Therefore
\[
 T\le8\pi X^{5/4}\frac{q(1+q)}{(1-q)^3}.
\]
Now
\[
 X=e^{2\pi}>e^6>2^6=64,
\]
so
\[
 q<e^{-192}<2^{-192}.
\]
Also $X=e^{2\pi}<e^8<3^8$ and $q<1/2$, whence
\[
 \frac{1+q}{(1-q)^3}<12.
\]
Thus
\[
 T<32\cdot3^{10}\cdot12\cdot2^{-192}.
\]
The exact integer inequality
\[
 32\cdot3^{10}\cdot12<2^{25}
\]
yields
\begin{equation}
 \boxed{T<2^{-167}}.
 \label{eq:g017-tail}
\end{equation}

\section{The fixed-point exclusion theorem}

On $[1/2,\pi]$ both $\Phi(y)$ and $\cos(y/2)$ are nonnegative.  Every possibly negative contribution therefore lies beyond $\pi$, and its absolute magnitude is bounded by $T$.  Combining \eqref{eq:g017-positive-block} and \eqref{eq:g017-tail},
\[
 F\!\left(-\frac14\right)
 \ge A-T
 >2^{-15}-2^{-167}
 >0.
\]
Hence
\begin{theorem}[SOH-G017 fixed-point exclusion]
\label{thm:g017-fixed-point-exclusion}
The quotient negative inversion $J(w)=1/(16w)$ has no fixed point in the root set of $F$.  Explicitly,
\[
 F(1/4)>0,
 \qquad
 \boxed{F(-1/4)>0}.
\]
The second inequality follows from the positive Riemann kernel and elementary analytic bounds, without using any list or ordinate of zeta zeros.
\end{theorem}

\section{Exceptional-orbit divisibility}

Recall the finite paired quotient set
\[
 P_J=\{w:F(w)=0,\;F(J(w))=0\}.
\]
Since $J^2=\mathrm{id}$, the involution $J$ acts on $P_J$.  By Theorem~\ref{thm:g017-fixed-point-exclusion}, this action is free.  Therefore $P_J$ is a disjoint union of two-cycles and
\begin{equation}
 \boxed{|P_J|\equiv0\pmod2}.
 \label{eq:g017-pj-even}
\end{equation}

Chapter~\ref{chap:g016-paired-spectrum-quotient} proves that
\[
 q(s)=\left(s-\frac12\right)^2
\]
maps the paired xi set $P_N$ exactly two-to-one onto $P_J$.  Hence
\[
 |P_N|=2|P_J|,
\]
and therefore
\begin{equation}
 \boxed{|P_N|\equiv0\pmod4}.
 \label{eq:g017-pn-four}
\end{equation}

More precisely, each quotient two-cycle $\{w,J(w)\}$ lifts to four distinct points in the $s$-plane.  Reflection $s\mapsto1-s$ exchanges the two points in each quotient fiber, while the canonical negative inversion exchanges the two fibers.  Thus every exceptional quotient two-cycle lifts to one full four-point orbit of the projective Klein-four action established in Chapter~\ref{chap:euler-riemann-negative-inversion-factorization}.

\section{Boundary of the result}

SOH-G017 excludes fixed exceptional roots and sharpens the structure of any remaining paired set.  It does not prove that $P_J$ or $P_N$ is empty.  It does not establish real-rootedness of $F$, PF$\infty$, SOH-G003, or the Riemann Hypothesis.
